\documentclass[11pt]{article}

\usepackage[letterpaper, margin=0.8in]{geometry}
\usepackage[T1]{fontenc}
\usepackage{charter}
\usepackage{microtype}
\usepackage{booktabs}
\usepackage{tabularx}
\usepackage{enumitem}
\usepackage{titlesec}
\usepackage{graphicx}
\usepackage{xcolor}
\usepackage{float}
\usepackage[hidelinks]{hyperref}

\definecolor{linkblue}{HTML}{1F5FA8}
\definecolor{changebg}{HTML}{FFF4D6}
\hypersetup{colorlinks=true, urlcolor=linkblue, linkcolor=linkblue}

\titleformat{\section}{\large\bfseries}{\thesection}{0.6em}{}
\titlespacing*{\section}{0pt}{12pt}{4pt}
\titleformat{\subsection}{\normalsize\bfseries}{\thesubsection}{0.6em}{}
\titlespacing*{\subsection}{0pt}{10pt}{2pt}
\setlength{\parindent}{0pt}
\setlength{\parskip}{5pt}
\setlist{nosep, leftmargin=1.4em, topsep=3pt}

% Highlight for anything that changed since Milestone 1
\newcommand{\changed}[1]{\colorbox{changebg}{\parbox{\dimexpr\linewidth-2\fboxsep}{\small\textbf{Changed since M1:} #1}}}

% ------------------------------------------------------------------
% Numbers and tables are written by scripts/build.py into tables/.
\newcommand{\DataMonths}{\textit{[run build.py]}}
\newcommand{\TripsRaw}{\textit{[n]}}
\newcommand{\TripsClean}{\textit{[n]}}
\newcommand{\NoEndStationPct}{\textit{[x]}}
\newcommand{\NoEndStationEbikePct}{\textit{[x]}}
\newcommand{\NStations}{\textit{[n]}}
\newcommand{\NCapMissing}{\textit{[n]}}
\newcommand{\NWeekdays}{\textit{[n]}}
\newcommand{\NWeekends}{\textit{[n]}}
\newcommand{\NoEndStationN}{\textit{[n]}}
\newcommand{\NoEndNoCoordsN}{\textit{[n]}}
\newcommand{\NoEndWithCoordsN}{\textit{[n]}}
\newcommand{\NoStartStationN}{\textit{[n]}}
\newcommand{\NOutsideNYC}{\textit{[n]}}
\newcommand{\NCapMissingNYC}{\textit{[n]}}
\newcommand{\NCapZero}{\textit{[n]}}
\newcommand{\NCapNotInFeed}{\textit{[n]}}
\IfFileExists{tables/meta.tex}{\renewcommand{\DataMonths}{January 2025, July 2025, October 2025 (3 months)}
\renewcommand{\TripsRaw}{11,844,428}
\renewcommand{\TripsClean}{11,826,946}
\renewcommand{\NoEndStationPct}{0.32}
\renewcommand{\NoEndStationEbikePct}{94}
\renewcommand{\NStations}{2,378}
\renewcommand{\NCapMissing}{70}
\renewcommand{\NWeekdays}{69}
\renewcommand{\NWeekends}{24}
\renewcommand{\NoEndStationN}{37,704}
\renewcommand{\NoEndNoCoordsN}{37,666}
\renewcommand{\NoEndWithCoordsN}{38}
\renewcommand{\NoStartStationN}{6,704}
\renewcommand{\NOutsideNYC}{68}
\renewcommand{\NCapMissingNYC}{67}
\renewcommand{\NCapZero}{26}
\renewcommand{\NCapNotInFeed}{41}
}{}

\newcommand{\datatable}[1]{%
  \IfFileExists{tables/#1.tex}{{\small\input{tables/#1.tex}}}%
  {\fbox{\parbox{0.9\linewidth}{\centering\small\textit{Run scripts/build.py to fill in tables/\detokenize{#1}.tex}}}}}

% Sketches live in ../sketches as q1.jpg / q1.png / q1.pdf (same for q2 ... q5)
\newcommand{\sketch}[2]{%
  \begin{figure}[H]\centering
  \IfFileExists{../sketches/#1.jpg}{\includegraphics[width=0.9\linewidth]{../sketches/#1.jpg}}{%
  \IfFileExists{../sketches/#1.png}{\includegraphics[width=0.9\linewidth]{../sketches/#1.png}}{%
  \IfFileExists{../sketches/#1.pdf}{\includegraphics[width=0.9\linewidth]{../sketches/#1.pdf}}{%
  \fbox{\parbox[c][1.8in][c]{0.8\linewidth}{\centering\small\textit{Sketch goes here: sketches/\detokenize{#1}.jpg}}}}}}
  \caption{#2}
  \end{figure}}
% ------------------------------------------------------------------

\begin{document}

\begin{center}
  {\Large\bfseries Spokes \& Flows: Milestone 2, Data Analysis and Sketches}\\[3pt]
  Team 2 \quad|\quad Harshith Kori Raj (hk4488), Lakshmi Hukunda Raju (lh41410)\\
  CS-GY 6313 B Information Visualization, Fall 2026\\[2pt]
  Code: \href{https://github.com/HarshithKoriRaj/spokesnflows}{github.com/HarshithKoriRaj/spokesnflows}
  (private, shared with \texttt{ctsilva} and \texttt{kimryan0416})
\end{center}

\section{What changed since Milestone 1}

The feedback said our plan was too big for two people and that our best line was ``where does the system
struggle, and when?'' We used that sentence to decide what to keep. Every change is also marked with a yellow
\emph{Changed since M1} box next to the question it affects.

\textbf{Audience.} Our main reader is someone planning the Citi Bike network, such as NYC DOT staff or a community
board deciding where to add docks. They need to see which stations keep losing or gaining bikes through riders'
trips, at what hours, and where those trips come from. A second reader is a regular rider who wants to know which
stations tend to drain at their commute time.

\textbf{Questions we kept, merged, refined and dropped.}
\begin{itemize}
  \item \textbf{Q1. When do people ride?} Kept. It sets up the daily rhythm the other views rely on.
  \item \textbf{Q2. Where and when do stations gain or lose bikes through trips?} Our old Q2 (where trips happen)
        and Q3 (imbalance) are now one map, as the feedback suggested.
  \item \textbf{Q3. Where do the trips behind that imbalance come from?} Our old Q4. We planned to focus on East
        River crossings, but the data showed they are only about 6.5\% of weekday morning trips
        (Table~\ref{tab:q3shares}). Most trips go to a nearby neighborhood in the same borough. So the question is
        now about neighborhood-to-neighborhood flows in general, and East River crossings are just one filter.
  \item \textbf{Q4. How much do weather and season change ridership?} Our old Q6. Kept because it is cheap (one row
        per day) and it tells a planner how much of a busy or quiet day is just the weather.
  \item \textbf{Q5 (optional). Do e-bike trips differ from classic-bike trips?} Our old Q5. The data part is done
        because it comes from the same cleaned table. We will only build the view if time allows.
  \item \textbf{Dropped:} our old Q7 (whether new stations connect to the rest of the network). It needs a second
        year of data and a network layout, and it doesn't help answer the main question. We also did not take on the
        optional congestion-pricing extension.
\end{itemize}

\textbf{Three months instead of a full year.} We use January, July and October 2025: one winter, one summer and one
fall month, \TripsRaw\ trips in total. This covers the seasonal range for Q4 while keeping the processing light.
Our pipeline works on any set of months, so the code would not change if we added more.

\section{Data handling}

\subsection{Sources}
\begin{itemize}
  \item \href{https://citibikenyc.com/system-data}{Citi Bike trip histories} published by Lyft. Monthly files from the
        \href{https://s3.amazonaws.com/tripdata/index.html}{public S3 bucket}: \texttt{202501}, \texttt{202507} and
        \texttt{202510-citibike-tripdata.zip}. One row is one trip.
  \item Dock capacity from the \href{https://gbfs.citibikenyc.com/gbfs/2.3/gbfs.json}{Citi Bike GBFS feed}
        (\texttt{station\_information}, downloaded November 2026).
  \item \href{https://data.cityofnewyork.us/d/9nt8-h7nd}{2020 Neighborhood Tabulation Areas} from NYC Open Data
        (262 areas, with borough names).
  \item \href{https://www.ncei.noaa.gov/cdo-web/datasets/GHCND/stations/GHCND:USW00094728/detail}{NOAA daily weather,
        Central Park}: daily high and low temperature, rain and snow.
\end{itemize}
How to get each file is listed in the repository README.

\subsection{Pipeline}
Everything runs in two scripts. \texttt{scripts/download.py} fetches the files. \texttt{scripts/build.py} converts
each month to Parquet, loads it into DuckDB, cleans it, joins the other datasets and writes the small tables the
visualizations will read. It also writes every table and number in this report. A full run on the three months
takes about 40 seconds on a laptop.

\begin{enumerate}
  \item \textbf{Filter:} keep trips that start in one of the three months (each file also holds a few trips that
        began the evening before), and drop trips longer than 3 hours, which are almost always bikes that weren't
        docked properly. The operator already removes trips under 60 seconds.
  \item \textbf{Derive:} trip duration, straight-line distance, start and end hour, day of week, and weekday/weekend.
  \item \textbf{Stations:} key each station on its ID, take its most common name, and use the median coordinates of
        classic-bike trips (classic bikes report the dock's position; e-bikes report their own GPS).
  \item \textbf{Join:} dock capacity from GBFS (by station ID, then by location within 75 m), neighborhood and
        borough by point-in-polygon, and weather by date.
  \item \textbf{Aggregate:} write the processed tables in Table~\ref{tab:map}.
\end{enumerate}

\begin{table}[H]\centering\small
\begin{tabular}{llrl}
\toprule
\textbf{Processed table} & \textbf{One row is} & \textbf{Rows} & \textbf{Used by} \\
\midrule
\texttt{hourly\_profile.csv} & day of week $\times$ hour $\times$ rider type & 336 & Q1 \\
\texttt{station\_hour.csv} & station $\times$ weekday/weekend $\times$ hour & 114,144 & Q2 \\
\texttt{stations.csv} & station (location, docks, neighborhood) & 2,378 & Q2, Q3 \\
\texttt{od\_nta.csv} & neighborhood pair $\times$ time period & 32,054 & Q3 \\
\texttt{daily.csv} & day, with weather & 93 & Q4 \\
\texttt{bike\_type\_summary.csv} & bike type $\times$ rider type & 4 & Q5 \\
\texttt{distance\_hist.csv} & bike type $\times$ 0.25 km distance bin & 66 & Q5 \\
\bottomrule
\end{tabular}
\caption{Every view reads from one of these tables, and all of them come from the same cleaned trip table.}\label{tab:map}
\end{table}

\subsection{What we found when checking the data}

\begin{table}[H]\centering
\datatable{quality}
\caption{Each check, how many trips it affects, and what we do. Percentages are of the raw trip count.}\label{tab:quality}
\end{table}

\textbf{Trips with no end station (feedback point 3).} \NoEndStationN\ trips (\NoEndStationPct\%) have no end
station, and \NoEndStationEbikePct\% of them are e-bikes. We expected these to be e-bikes left outside a dock, but
\NoEndNoCoordsN\ of them also have no end coordinates, so we can't tell where they ended. Only \NoEndWithCoordsN\
trips end at a known location away from a station. The classic-bike trips in this group last about 25 hours, which
looks like the system closing trips for bikes that were never returned. So the 2025 files have essentially no
``dockless'' trips with a known end point. We keep these trips in the time and weather views (Q1, Q4) and leave them
out of the station and neighborhood views (Q2, Q3), which slightly undercounts arrivals there. Likewise,
\NoStartStationN\ trips have no start station and no start coordinates. They are handled the same way.

\textbf{Stations outside NYC.} \NOutsideNYC\ stations fall outside every NYC neighborhood. They are Jersey City and
Hoboken docks that some NYC trips ended at. They stay out of the map, and those trips count as leaving NYC.

\textbf{Dock capacity.} The GBFS feed lists today's stations, not 2025's. \NCapMissingNYC\ NYC stations still have no
capacity after matching. \NCapZero\ of them are listed with zero docks (probably closed or temporary), and
\NCapNotInFeed\ are no longer in the feed. A few are busy, such as West St \& Chambers St. These stations stay on the map in grey,
labeled ``no dock data'', and get no color.

\subsection{What our imbalance measure shows, and what it doesn't (feedback point 2)}

For each station and hour we compute \textbf{net trip flow} = arrivals $-$ departures, averaged over the
\NWeekdays\ weekdays (or \NWeekends\ weekend days), and divide by the number of docks. Adding it up from midnight gives
the \textbf{cumulative net trips per dock} at each hour. A value of $-0.8$ at 10:00 means that, on an average
weekday, riders took 0.8 more bikes per dock out of the station than they brought back since midnight.

This measures \emph{trips}, not how many bikes were at the station. Our own data shows the gap clearly. By 10:00 on
weekdays, North Moore St \& Greenwich St in Tribeca gains 3.65 bikes per dock through trips
(Table~\ref{tab:q2extremes}). That is far more bikes than the station can hold, so staff must be taking bikes away
(by rebalancing or valet service). Rebalancing trips are not in the public data. So we will say a station
``loses (or gains) bikes through riders' trips'', never that it is ``empty'' or ``full'', and the map will carry a note
saying so.

\section{Domain questions: processed data and sketches}

For each question we give the unit of observation, the transformation, a sample of the processed table, and our
sketch. Across all views: panels that are compared share one scale, missing data is drawn as its own grey category,
and every color difference is also shown another way (shape, line style or position).

\subsection*{Q1. When do people ride, and how do members and casual riders differ?}
\textbf{Unit:} day of week $\times$ hour $\times$ rider type.
\textbf{Transformation:} count trips by start hour and weekday, then divide by how many times that weekday appears in
the three months to get trips per day.
\textbf{What the sample shows:} on weekdays members have two peaks, about 9,200 trips in the 8:00 hour and 11,000 at
17:00. Casual riders barely have a morning peak (about 900 at 8:00) but rise in the evening.
\textbf{View:} two heatmaps side by side (members, casual riders), hour across and day of week down. Color is each
hour's share of that group's weekly trips, so the shapes can be compared even though members take 83\% of all trips.
Panel titles give trips per day. Hovering a cell shows the exact count.

\begin{table}[H]\centering
\datatable{q1_sample}
\caption{Sample of \texttt{hourly\_profile.csv} (Wednesday and Saturday at 8:00 and 17:00). The full table has 336 rows.}
\end{table}
\sketch{q1}{Q1 sketch: hour $\times$ weekday heatmaps for members and casual riders.}

\subsection*{Q2. Where and when do stations gain or lose bikes through trips?}
\changed{merges old Q2 and Q3 into one map; the measure is now explicitly a trip measure, normalized by docks.}

\textbf{Unit:} station $\times$ weekday/weekend $\times$ hour (2,378 stations $\times$ 48).
\textbf{Transformation:} departures are counted at the start station and start hour, and arrivals at the end
station and end hour. Both are divided by the number of days. Net and cumulative net are then divided by the number
of docks. Every station has all 24 hours, so hours with no trips appear as zero instead of disappearing.
\textbf{What the sample shows:} the busiest station, W 21 St \& 6 Ave (74 docks), gains about 0.66 bikes per dock
through trips by noon on weekdays. Across all stations with dock data, 91 lose more than half a bike per dock by
10:00 and 79 gain more than that. The top losers are in residential areas such as the Upper East Side and Greenpoint.
\textbf{View:} a map of stations colored on a diverging scale centered at zero (red = losing, blue = gaining, legend
in net trips per dock). Losing stations are drawn as downward triangles and gaining ones as upward triangles, so
direction doesn't depend on color. Stations with no dock data are grey hollow circles. An hour slider and a
weekday/weekend toggle sit above the map. Clicking a station opens a bar chart of its departures and arrivals by
hour, with a line for cumulative net per dock.

\begin{table}[H]\centering
\datatable{q2_sample}
\caption{Sample of \texttt{station\_hour.csv}: the busiest station, weekday mornings (per average weekday).}
\end{table}

\begin{table}[H]\centering
\datatable{q2_extremes}
\caption{Stations with at least 15 docks that lose or gain the most bikes per dock through trips by 10:00 on an
average weekday.}\label{tab:q2extremes}
\end{table}
\sketch{q2}{Q2 sketch: net-flow map with hour slider and weekday/weekend toggle, and the hourly profile shown when a station is selected.}

\subsection*{Q3. Where do the trips behind that imbalance come from?}
\changed{we dropped the East River crossings as the focus because they are only about 6.5\% of weekday morning trips
(Table~\ref{tab:q3shares}). Crossings are now a filter, and the question covers all neighborhood-to-neighborhood flows.}

\textbf{Unit:} start neighborhood $\times$ end neighborhood $\times$ time period (weekday 7--10, weekday 16--19, other
weekday hours, weekend).
\textbf{Transformation:} join each trip's start and end station to its 2020 neighborhood and borough, count trips per
pair and period, divide by the number of days, and flag pairs that cross a borough line or the East River.
\textbf{What the sample shows:} about 70\% of trips go to a different neighborhood in the same borough and 21\% stay in
the same neighborhood. The busiest morning flows are short hops into Midtown, from Chelsea-Hudson Yards, Hell's
Kitchen and Murray Hill. On weekday mornings Brooklyn sends about 1,040 trips per day into Manhattan, against about
15,200 trips within Manhattan.
\textbf{View:} a borough-to-borough matrix on the left and a map of the busiest neighborhood flows on the right. Flows
are curved lines whose width is trips per day. East River crossings are solid and other lines dashed. A period
selector changes both views, and clicking a neighborhood keeps only its flows. Selecting a station in Q2 highlights
its neighborhood here.

\begin{table}[H]\centering
\datatable{q3_shares}
\caption{Where trips go, by period (trips with a known start and end neighborhood). ``Same borough'' means a
different neighborhood in the same borough.}\label{tab:q3shares}
\end{table}

\begin{table}[H]\centering
\datatable{q3_sample}
\caption{Sample of \texttt{od\_nta.csv}: the busiest pairs of different neighborhoods, weekday 7--10.}
\end{table}

\begin{table}[H]\centering
\datatable{q3_boroughs}
\caption{Weekday 7--10 trips per day between boroughs (rows: start, columns: end).}
\end{table}
\sketch{q3}{Q3 sketch: borough matrix and neighborhood flow map with period selector.}

\subsection*{Q4. How much do weather and season change ridership?}
\changed{uses three months (January, July, October) instead of the full year.}

\textbf{Unit:} one day (93 days).
\textbf{Transformation:} count trips per day by rider type and bike type, then join that day's high temperature, rain
and snow from Central Park by date. Holidays are flagged so they aren't mistaken for weather effects.
\textbf{What the data shows:} an average day had about 68,500 trips in January, 160,600 in July and 152,500 in October.
On rainy weekdays (0.1 inch or more), members rode at 88\% of their usual level for that month, but casual riders
dropped to 68\%.
\textbf{View:} a calendar heatmap of daily trips for the three months (weeks across, days of the week down, holidays
outlined). Below it, two scatter plots share one y-axis: members and casual riders, each showing trips as a percent
of that group's monthly average against the day's high temperature. Rainy days are filled dots and dry days hollow.

\begin{table}[H]\centering
\datatable{q4_sample}
\caption{Sample of \texttt{daily.csv} (first ten days of January; Jan 1 is a holiday).}
\end{table}
\sketch{q4}{Q4 sketch: calendar heatmap and trips vs.\ temperature for members and casual riders.}

\subsection*{Q5 (optional). Do e-bike trips differ from classic-bike trips?}
\textbf{Unit:} bike type $\times$ rider type (summary), and bike type $\times$ 0.25 km distance bin (histogram).
\textbf{Transformation:} straight-line distance between start and end coordinates, which is shorter than the real route.
\textbf{What the data shows:} e-bikes carry about 71\% of trips. For members, the median e-bike trip covers 1.63 km
against 1.06 km on a classic bike, and takes only about one minute longer.
\textbf{View:} overlaid step histograms of trip distance, solid for e-bikes and dashed for classic bikes, with the
medians marked.

\begin{table}[H]\centering
\datatable{q5_sample}
\caption{\texttt{bike\_type\_summary.csv}. ``No end station'' is the share of trips without an end station.}
\end{table}
\sketch{q5}{Q5 sketch (optional): trip distance by bike type.}

\section*{AI use}
{\small
We used Claude (Anthropic) to help write the two data scripts in our repository's \texttt{scripts/} folder
(\texttt{download.py} and \texttt{build.py}), to run the first data checks, and to draft parts of this report's
text. We reviewed and ran the code, checked the outputs against the raw files, and edited the text, and we can
explain every step. The sketches are our own drawings, made without AI. The same statement is in the repository's
README.
}

\end{document}
